\src{254}
Mặt khác, nhắc lại rằng nếu $\mathcal F$ là một bó coherent trên $X$, thì nó
đẳng cấu với một bó thương của bó $\mathcal O_X(-n)^N$, với $n$ và $N$ đủ
lớn. Do đó, các bó $\mathcal F_s$ cảm sinh trên các thớ cũng là thương của bó
$\mathcal O(-n)^N$ trên thớ.

Từ hai nhận xét này, ta suy ra chiều «chỉ khi» của định lý sau:

\result{ĐỊNH LÝ}{2.1} Cho $X$ xạ ảnh trên $S$ Noether, với
$\mathcal O_X(1)$ rất rộng trên $X$ tương đối với $S$. Cho $E$ là một tập hợp
các lớp bó trên các thớ của $X/S$. Để $E$ bị chặn, điều kiện cần và đủ\footnote{Văn
bản in chỉ ghi «điều kiện cần»; bản đính chính năm 1962 cho Exposé 221 quy định
«điều kiện cần và đủ».} là nó thỏa mãn các điều kiện sau:

\medskip
\noindent (a) Tồn tại một bó coherent $\mathcal L$ trên $X$ (có thể giả sử có
dạng $\mathcal O_X(-n)^N$) sao cho $E$ được chứa trong họ các lớp của những bó
coherent là thương của các bó dạng $\mathcal L_K$.

\smallskip
\noindent (b) Các đa thức Hilbert $P_{\mathcal F}$ của những bó $\mathcal F$
có lớp thuộc $E$ đều thuộc cùng một tập hợp hữu hạn các đa thức.

Còn phải chứng minh chiều «nếu», đây sẽ là một trường hợp riêng của một kết quả
chính xác hơn. Với mọi môđun coherent $\mathcal F$ trên một tiền lược đồ thuộc
kiểu hữu hạn trên một trường $K$, và mọi số nguyên $r$, gọi $N_r$ là môđun con
của $\mathcal F$ mà các tiết diện trên một tập mở là những tiết diện của
$\mathcal F$ trên tập mở đó có giá mang chiều $<r$. Khi đó
$N_r=\mathcal F$ nếu $r>\dim\operatorname{supp}\mathcal F$, còn $N_r=0$ nếu
$r\leq0$, và ta thu được một lọc tăng hữu hạn của $\mathcal F$ mà các thương
$N_{r+1}/N_r$\footnote{Văn bản in đảo thương thành $N_r/N_{r+1}$, điều này
không tương thích với lọc tăng vừa được định nghĩa.} có các chu trình nguyên tố
liên kết đúng bằng những chu trình nguyên tố liên kết với $\mathcal F$ mang
chiều $r$. Ta đặt
\[
\mathcal F_{(r)}=\mathcal F/N_r,
\]
thì các chu trình nguyên tố liên kết với $\mathcal F_{(r)}$ đúng bằng những chu
trình nguyên tố liên kết với $\mathcal F$ mang chiều $\geq r$; đặc biệt, nó
bằng $\mathcal F$ khi và chỉ khi các chu trình nguyên tố liên kết với
$\mathcal F$ mang chiều $\geq r$. Với ký hiệu này:

\result{ĐỊNH LÝ}{2.2} Dưới các giả thiết của (2.1), cho $s$ là một số nguyên
và giả sử $E$ thỏa mãn điều kiện (a) cùng dạng yếu hơn sau đây của (b):

\medskip
\noindent $(b_s)$ Các đa thức Poincaré $P_{\mathcal F}$ của những bó
$\mathcal F$ có lớp thuộc $E$ có các hệ số ở bậc $\geq s-1$ được chặn.

Dưới các điều kiện này, các bó $\mathcal F_{(s)}$ (khi lớp của $\mathcal F$
chạy trong $E$) tạo thành một họ bị chặn. Hơn nữa, các hệ số bậc $s-2$ của
các $P_{\mathcal F}$ bị chặn dưới.
